% ==========================================================================
% Relatorio Tecnico - Padrao UFAL
% Framework Whitelabel para Sistemas Online de Aluguel de Casas e Apartamentos
% Disciplina: Reuso de Software e Metodologias Ageis - 2026.1
% Alunos: Marcos Melo dos Santos e Leonardo Barbosa
%
% Compilacao: pdflatex relatorio.tex (rodar 2x para sumario/referencias)
% ==========================================================================
\documentclass[12pt,a4paper]{article}

% ----------------------------- Pacotes ------------------------------------
\usepackage[utf8]{inputenc}
\usepackage[T1]{fontenc}
\usepackage[brazil]{babel}
\usepackage[a4paper,margin=2.5cm]{geometry}
\usepackage{graphicx}
\usepackage{xcolor}
\usepackage{hyperref}
\usepackage{listings}
\usepackage{enumitem}
\usepackage{booktabs}
\usepackage{longtable}
\usepackage{fancyhdr}
\usepackage{titlesec}
\usepackage{caption}

% --------------------------- Hyperlinks -----------------------------------
\hypersetup{
    colorlinks=true,
    linkcolor=black,
    citecolor=black,
    urlcolor=blue,
    pdftitle={Framework Whitelabel de Aluguel - UFAL 2026.1},
    pdfauthor={Marcos Melo dos Santos e Leonardo Barbosa}
}

% ------------------------- Estilo de codigo -------------------------------
\definecolor{codegray}{rgb}{0.95,0.95,0.95}
\definecolor{codegreen}{rgb}{0,0.5,0}
\definecolor{codeblue}{rgb}{0,0,0.7}
\lstset{
    backgroundcolor=\color{codegray},
    basicstyle=\ttfamily\footnotesize,
    keywordstyle=\color{codeblue}\bfseries,
    commentstyle=\color{codegreen}\itshape,
    breaklines=true,
    frame=single,
    showstringspaces=false,
    tabsize=2,
    captionpos=b
}

% ------------------------- Cabecalho/Rodape -------------------------------
\pagestyle{fancy}
\fancyhf{}
\lhead{\small Framework Whitelabel de Aluguel}
\rhead{\small UFAL 2026.1}
\cfoot{\thepage}

% ==========================================================================
\begin{document}

% ----------------------------- Capa ---------------------------------------
\begin{titlepage}
    \centering
    \vspace*{1cm}
    {\Large \textbf{UNIVERSIDADE FEDERAL DE ALAGOAS}}\\[0.3cm]
    {\large Instituto de Computacao}\\[2.5cm]

    {\Huge \textbf{Framework Whitelabel para Sistemas Online de Aluguel de Casas e Apartamentos}}\\[0.5cm]
    {\large Relatorio Tecnico de Desenvolvimento}\\[3cm]

    \begin{flushright}
        \begin{minipage}{0.6\textwidth}
            \textbf{Alunos:}\\
            Marcos Melo dos Santos\\
            {\small (Gateway, Reservas Core e Front-end via Figma/MCP)}\\
            Leonardo Barbosa\\
            {\small (Persistencia e logica interna de Auth e Catalogo)}\\[1cm]
            \textbf{Disciplina:}\\
            Reuso de Software e Metodologias Ageis\\[1cm]
            \textbf{Periodo:} 2026.1
        \end{minipage}
    \end{flushright}

    \vfill
    {\large Maceio -- Alagoas\\ 2026}
\end{titlepage}

% ---------------------------- Sumario -------------------------------------
\tableofcontents
\newpage

% ==========================================================================
\section{Enunciado do Problema}
% ==========================================================================
O presente trabalho tem como objetivo o desenvolvimento de um \textbf{Framework
Whitelabel} para a construcao de sistemas online de aluguel de casas e apartamentos.
Um produto \textit{whitelabel} e aquele que pode ser reaproveitado e customizado por
diferentes clientes (imobiliarias, plataformas de temporada, redes de coliving),
mantendo um nucleo comum de regras de negocio.

O desafio central, do ponto de vista de \textbf{Reuso de Software}, e projetar uma
arquitetura que maximize o reaproveitamento do fluxo de reserva, ao mesmo tempo em que
permita variacoes especificas de cada modalidade de locacao. Para isso, adota-se o
padrao de projeto \textit{Template Method}, que define o esqueleto invariante do
algoritmo de reserva e delega passos variaveis (\textit{hotspots}) para subclasses.

O sistema e estruturado como um \textbf{monorepo} de microsservicos em PHP 8.2 e MySQL,
contemplando autenticacao, catalogo de imoveis e a engine de reservas, integrados por
um \textit{API Gateway}.

% ==========================================================================
\section{Planejamento Scrum por Prioridades}
% ==========================================================================
O desenvolvimento seguiu uma abordagem agil incremental, organizada em um
\textit{Product Backlog} priorizado. Cada item foi tratado como uma \textit{Sprint}
curta, com entrega versionada (branch + commit) ao final de cada etapa.

\begin{longtable}{c p{7.5cm} c}
    \toprule
    \textbf{Prioridade} & \textbf{Item do Backlog (User Story)} & \textbf{Status} \\
    \midrule
    \endhead
    1 (Alta)  & Inicializacao do repositorio: \texttt{.gitignore}, \texttt{README} e relatorio tecnico. & Concluido \\
    2 (Alta)  & API Gateway: roteamento central \textit{stateless} e comunicacao interna via cURL. & Concluido \\
    3 (Alta)  & Microsservico de Autenticacao (\texttt{ms-auth}) com JWT e banco \texttt{db\_auth}. & Concluido \\
    4 (Media) & Microsservico de Catalogo (\texttt{ms-catalogo}) e banco \texttt{db\_catalogo}. & Concluido \\
    5 (Alta)  & Engine de Reservas (\texttt{ms-reservas}): classe abstrata \texttt{RentEngine} (Template Method). & Concluido \\
    6 (Media) & Classes concretas \texttt{LongTermRent} e \texttt{VacationRent}, provando o reuso. & Concluido \\
    7 (Baixa) & Front-end integrado no Gateway (Blade/HTML) consumindo o ecossistema. & Concluido \\
    \bottomrule
    \caption{Product Backlog priorizado do framework.}
\end{longtable}

\subsection{Justificativa das Prioridades}
Os itens de infraestrutura e seguranca (Gateway e Autenticacao) recebem prioridade
alta por serem dependencias transversais. A engine de reservas, sendo o nucleo de
reuso e o diferencial academico do trabalho, tambem possui prioridade alta. O catalogo
e as classes concretas dependem dessas bases, e o front-end e a camada de apresentacao
final, com prioridade mais baixa por nao bloquear as demais entregas.

% ==========================================================================
\section{Tecnicas de Reuso de Software Aplicadas}
% ==========================================================================
O projeto aplica multiplas tecnicas de reuso, em diferentes niveis de granularidade:

\begin{description}[leftmargin=1.5cm]
    \item[Frameworks de Aplicacao] O proprio produto e um \textit{framework whitelabel}:
    fornece um esqueleto reutilizavel que e instanciado por diferentes aplicacoes de aluguel.
    \item[Padroes de Projeto] Uso do \textit{Template Method} para fixar o fluxo de
    reserva e variar passos especificos via \textit{hotspots}.
    \item[Reuso por Heranca] As modalidades \texttt{VacationRent} e \texttt{LongTermRent}
    herdam de \texttt{RentEngine}, reaproveitando o algoritmo comum.
    \item[Reuso por Composicao] A funcionalidade \texttt{CreditCheckComponent} e injetada
    em \texttt{LongTermRent} por composicao, evitando heranca rigida e favorecendo flexibilidade.
    \item[Reuso de Servicos (SOA)] A arquitetura de microsservicos permite que
    autenticacao e catalogo sejam reutilizados por qualquer cliente do ecossistema.
\end{description}

% ==========================================================================
\section{Arquitetura da Solucao}
% ==========================================================================
A solucao adota uma arquitetura de \textbf{microsservicos} em \textit{monorepo},
com isolamento de dados (\textit{Database per Service}). O cliente comunica-se
exclusivamente com o Gateway, que orquestra as chamadas internas.

\begin{table}[h]
    \centering
    \begin{tabular}{l c l l}
        \toprule
        \textbf{Servico} & \textbf{Porta} & \textbf{Banco} & \textbf{Responsabilidade} \\
        \midrule
        \texttt{gateway}     & 8003 & ---            & Roteamento e composicao (cURL) \\
        \texttt{ms-auth}     & 8000 & \texttt{db\_auth}     & Autenticacao JWT \\
        \texttt{ms-catalogo} & 8001 & \texttt{db\_catalogo} & Catalogo de imoveis \\
        \texttt{ms-reservas} & 8002 & \texttt{db\_reservas} & Engine de reservas \\
        \bottomrule
    \end{tabular}
    \caption{Mapeamento de servicos, portas e bancos de dados.}
\end{table}

\begin{lstlisting}[caption={Topologia logica de comunicacao.},label={lst:topologia}]
[ Cliente / Browser ]
         |
         v
[ API Gateway :8003 ]  --- stateless, roteia via cURL
   |        |        |
   v        v        v
ms-auth  ms-catalogo  ms-reservas
 :8000      :8001        :8002
db_auth   db_catalogo  db_reservas
\end{lstlisting}

% ==========================================================================
\section{API Gateway}
% ==========================================================================
O \textit{API Gateway} e o unico ponto de entrada do sistema (\textit{Single Entry
Point}). Suas caracteristicas principais sao:

\begin{itemize}
    \item \textbf{Stateless:} nao mantem estado de sessao; cada requisicao carrega o
    token JWT necessario, repassado aos microsservicos.
    \item \textbf{Roteamento:} mapeia rotas externas (ex.: \texttt{/api/imoveis}) para
    os servicos internos correspondentes.
    \item \textbf{Comunicacao via cURL:} as chamadas internas sao feitas por HTTP com a
    extensao cURL do PHP nativo, mantendo baixo acoplamento.
    \item \textbf{Apresentacao:} serve as \textit{views} (Blade/HTML) do front-end
    integrado.
\end{itemize}

\subsection{Microsservico de Autenticacao (ms-auth)}
O \texttt{ms-auth} (porta 8000, banco \texttt{db\_auth}) e responsavel pela identidade
do ecossistema. Implementa JWT no algoritmo HS256 de forma autocontida, expondo os
endpoints \texttt{/register}, \texttt{/login}, \texttt{/me} e \texttt{/validate}. As
senhas sao armazenadas com \texttt{bcrypt} e o acesso a dados usa PDO com
\textit{prepared statements}. Por ser stateless, o token viaja em cada requisicao
(cabecalho \texttt{Authorization}), sendo repassado pelo Gateway aos demais servicos.

\subsection{Microsservico de Catalogo (ms-catalogo)}
O \texttt{ms-catalogo} (porta 8001, banco \texttt{db\_catalogo}) gerencia o cadastro de
imoveis (casas e apartamentos) com uma API REST de CRUD: \texttt{GET /} (com filtros por
cidade, tipo e disponibilidade), \texttt{GET /\{id\}}, \texttt{POST /}, \texttt{PUT
/\{id\}} e \texttt{DELETE /\{id\}}. O relacionamento com o usuario proprietario
(\texttt{owner\_id}) e logico, sem chave estrangeira fisica, respeitando o isolamento do
padrao \textit{Database per Service}.

% ==========================================================================
\section{Front-End Integrado}
% ==========================================================================
A camada de apresentacao e servida pelo proprio Gateway (porta 8003). O design foi
derivado via \textbf{Figma MCP} (\texttt{get\_figma\_data}), mesclando dois kits da
Figma Community com paleta de cores propria (identidade UFAL):

\begin{itemize}
    \item \textbf{M-Rent} (Property Management): cards de imovel com \textit{border-radius}
    12px, badges de tipo e estrutura de listagem com filtros.
    \item \textbf{Estatery} (Real Estate SaaS UI Kit): \textit{hero}, \textit{Search Bar},
    abas por tipo de imovel e barra de resultados.
\end{itemize}

A view consome exclusivamente as rotas \texttt{/api/*}:

\begin{itemize}
    \item \textbf{Autenticacao:} login via \texttt{POST /api/auth/login}; token JWT no
    \texttt{localStorage}, reenviado em \texttt{Authorization}.
    \item \textbf{Catalogo:} busca com filtros de API (\texttt{city}, \texttt{type},
    \texttt{available}) e filtros client-side (preco, quartos, area, ordenacao).
    \item \textbf{Reservas:} criacao via \texttt{POST /api/reservas} com modalidades
    \texttt{vacation} e \texttt{long\_term}.
\end{itemize}

Assets em \texttt{gateway/public/assets/css/app.css} e \texttt{gateway/views/home.php}.

% ==========================================================================
\section{Template Method e Hotspots}
% ==========================================================================
O nucleo de reuso e a classe abstrata \texttt{RentEngine}. Ela implementa o
\textbf{Template Method} \texttt{processReservation()}, que fixa a sequencia
invariante de passos do processo de reserva. Os \textbf{hotspots} (metodos abstratos
ou \textit{hooks}) representam os pontos de variacao implementados por cada subclasse.

\begin{description}[leftmargin=1.5cm]
    \item[Frozen Spots] Passos fixos: validacao de disponibilidade, persistencia da
    reserva e confirmacao. Sao implementados na classe abstrata e nao mudam.
    \item[Hot Spots] Passos variaveis: calculo de preco, regras de validacao
    especificas e etapas opcionais (ex.: analise de credito). Sao abstratos ou
    \textit{hooks}, definidos pelas subclasses.
\end{description}

\begin{lstlisting}[language=PHP,caption={Template Method implementado em \texttt{ms-reservas/src/Engine/RentEngine.php}.},label={lst:template}]
abstract class RentEngine
{
    // TEMPLATE METHOD: esqueleto invariante, final (nao sobrescrevivel)
    final public function processReservation(ReservationRequest $req): ReservationResult
    {
        $this->validateBaseRequest($req);          // (1) FROZEN
        $this->ensureAvailability($req);           // (2) FROZEN
        $price  = $this->calculatePrice($req);     // (3) HOTSPOT (abstrato)
        $this->applyBusinessRules($req, $price);   // (4) HOTSPOT (hook)
        $status = $this->resolveStatus($req, $price); // (5) HOTSPOT (hook)
        $id = $this->repository->create(           // (6) FROZEN: persistencia
            $req, $this->modalityName(), $price, $status, $this->reservationNotes($req)
        );
        $this->afterReservation($id, $req, $price); // (7) HOTSPOT (hook)
        return new ReservationResult(/* ... */);
    }

    // HOTSPOTS obrigatorios (abstratos)
    abstract protected function modalityName(): string;
    abstract protected function calculatePrice(ReservationRequest $req): float;

    // HOTSPOTS opcionais (hooks com comportamento default)
    protected function applyBusinessRules(ReservationRequest $req, float $price): void {}
    protected function resolveStatus(ReservationRequest $req, float $price): string { return 'confirmed'; }
    protected function afterReservation(int $id, ReservationRequest $req, float $price): void {}
}
\end{lstlisting}

A selecao da modalidade concreta ocorre em \texttt{EngineFactory}, ponto unico de
extensao do framework: registrar uma nova modalidade nao exige alterar o nucleo,
respeitando o \textit{Open/Closed Principle}.

\subsection{Classes Concretas e Prova de Reuso}
Duas modalidades concretas instanciam o framework, reaproveitando o mesmo Template
Method e variando apenas os hotspots:

\begin{description}[leftmargin=1.5cm]
    \item[\texttt{VacationRent} (temporada)] Reuso por \textbf{heranca}: sobrescreve
    somente os hotspots obrigatorios (\texttt{modalityName} e \texttt{calculatePrice}),
    aplicando adicional de alta temporada para estadias curtas e taxa de limpeza.
    Os hooks opcionais usam o comportamento default da \texttt{RentEngine}.
    \item[\texttt{LongTermRent} (longa duracao)] Combina \textbf{heranca} e
    \textbf{composicao}: calcula o preco mensal com desconto para contratos longos e
    injeta o \texttt{CreditCheckComponent} no hook \texttt{applyBusinessRules}. O
    resultado da analise de credito governa o hook \texttt{resolveStatus}
    (\texttt{confirmed} se aprovado, \texttt{pending} caso contrario).
\end{description}

A Tabela~\ref{tab:reuso} resume como cada modalidade exercita os pontos do framework.

\begin{table}[h]
    \centering
    \begin{tabular}{l c c}
        \toprule
        \textbf{Ponto do framework} & \textbf{VacationRent} & \textbf{LongTermRent} \\
        \midrule
        \texttt{calculatePrice} (abstract) & sobrescrito (diaria) & sobrescrito (mensal) \\
        \texttt{applyBusinessRules} (hook)  & default (vazio)      & credit check (composicao) \\
        \texttt{resolveStatus} (hook)       & default (confirmed)  & dinamico (credito) \\
        Esqueleto (frozen spots)            & reutilizado          & reutilizado \\
        \bottomrule
    \end{tabular}
    \caption{Reuso dos pontos do framework por modalidade.}
    \label{tab:reuso}
\end{table}

\begin{lstlisting}[language=PHP,caption={Reuso por composicao: credit check no hook do LongTermRent.},label={lst:longterm}]
final class LongTermRent extends RentEngine
{
    public function __construct(
        ReservationRepository $repository,
        private readonly CreditCheckComponent $creditCheck = new CreditCheckComponent()
    ) { parent::__construct($repository); }

    protected function applyBusinessRules(ReservationRequest $req, float $price): void {
        $this->creditResult = $this->creditCheck->evaluate($req); // COMPOSICAO
    }
    protected function resolveStatus(ReservationRequest $req, float $price): string {
        return $this->creditResult['approved'] ? 'confirmed' : 'pending';
    }
}
\end{lstlisting}

Na Etapa 6, \texttt{VacationRent} sobrescreve \texttt{calculatePrice()} com tarifas de
temporada, enquanto \texttt{LongTermRent} adiciona, por composicao, o
\texttt{CreditCheckComponent} dentro do \textit{hook} \texttt{applyBusinessRules()},
demonstrando a combinacao de reuso por heranca e por composicao.

% ==========================================================================
\section{Guia do Desenvolvedor}
% ==========================================================================
\subsection{Pre-requisitos}
PHP 8.2+ (extensoes \texttt{pdo\_mysql}, \texttt{curl}, \texttt{json}), MySQL 8.0+ e
Composer 2.x.

\subsection{Subindo os Servicos}
Em terminais separados:
\begin{lstlisting}[language=bash]
cd ms-auth      && composer install && php -S localhost:8000 -t public
cd ms-catalogo  && composer install && php -S localhost:8001 -t public
cd ms-reservas  && composer install && php -S localhost:8002 -t public
cd gateway      && composer install && php -S localhost:8003 -t public
\end{lstlisting}

\subsection{Criando os Bancos}
\begin{lstlisting}[language=bash]
mysql -u root -p < ms-auth/sql/db_auth.sql
mysql -u root -p < ms-catalogo/sql/db_catalogo.sql
mysql -u root -p < ms-reservas/sql/db_reservas.sql
\end{lstlisting}

\subsection{Como Estender o Framework}
Para criar uma nova modalidade de aluguel, basta criar uma subclasse de
\texttt{RentEngine} e implementar os \textit{hotspots}, sem alterar o fluxo principal,
respeitando o \textit{Open/Closed Principle}.

% ==========================================================================
\section{Divisao de Responsabilidades e Pendencias}
% ==========================================================================
O trabalho foi conduzido em dupla, com divisao clara de escopo. \textbf{Marcos Melo dos
Santos} ficou responsavel pelo \textit{API Gateway}, pelo nucleo de reuso (a engine
\texttt{RentEngine} com Template Method e suas modalidades concretas) e pelo front-end
integrado. \textbf{Leonardo Barbosa} ficou responsavel pela persistencia e pela logica
interna dos servicos de \textbf{Autenticacao} e \textbf{Catalogo} (banco de dados,
autorizacao e criptografia).

O ecossistema executa de ponta a ponta. Contudo, alguns pontos sob responsabilidade do
Leonardo ainda precisam ser concluidos ou validados para uma operacao real e segura. Vale
notar que o codigo de Auth e Catalogo \textit{ja existe} (rotas, PDO, \texttt{bcrypt} e
JWT HS256); os itens da Tabela~\ref{tab:pendencias} representam o que efetivamente resta.

\begin{table}[h]
    \centering
    \begin{tabular}{l p{8.5cm}}
        \toprule
        \textbf{Area} & \textbf{Pendencia} \\
        \midrule
        Config.     & Criar \texttt{.env} reais; definir \texttt{JWT\_SECRET} forte; provisionar/validar \texttt{db\_auth} e \texttt{db\_catalogo} em MySQL real. \\
        Auth        & Aplicar autorizacao por papel (\texttt{role}); seed do admin sem hash fixo; refresh/logout (revogacao); mapear \texttt{PDOException} para HTTP. \\
        Catalogo    & Proteger escrita (\texttt{POST/PUT/DELETE}) com token/role; vincular \texttt{owner\_id} ao usuario autenticado; paginacao/ordenacao na listagem. \\
        Integracao  & Acionar \texttt{POST /api/auth/validate} como \textit{middleware}: o Gateway hoje apenas repassa o cabecalho \texttt{Authorization}, sem validar o token. \\
        \bottomrule
    \end{tabular}
    \caption{Pendencias reais sob responsabilidade do Leonardo (Auth e Catalogo).}
    \label{tab:pendencias}
\end{table}

% ==========================================================================
\section{Referencias}
% ==========================================================================
\begin{thebibliography}{9}
\bibitem{gof} GAMMA, E. et al. \textit{Design Patterns: Elements of Reusable
Object-Oriented Software}. Addison-Wesley, 1994.
\bibitem{sommerville} SOMMERVILLE, I. \textit{Engenharia de Software}. 10. ed.
Pearson, 2018.
\bibitem{newman} NEWMAN, S. \textit{Building Microservices}. 2. ed. O'Reilly, 2021.
\bibitem{schwaber} SCHWABER, K.; SUTHERLAND, J. \textit{The Scrum Guide}. 2020.
\bibitem{fowler} FOWLER, M. \textit{Patterns of Enterprise Application Architecture}.
Addison-Wesley, 2002.
\bibitem{php} PHP DOCUMENTATION. \textit{PHP 8.2 Manual}. Disponivel em:
\url{https://www.php.net/manual/}. Acesso em: 2026.
\end{thebibliography}

\end{document}
